\documentclass[732]{studrep}

\RequirePackage{tabularx}
\RequirePackage{graphicx}

\graphicspath{{pics/}}

% \usepackage{showframe}

% pygmetize
\RequirePackage{minted}
\RequirePackage{indentfirst}
\usemintedstyle{tango} % bw

% \setmainfont{Times New Roman}

\setminted{autogobble,mathescape,linenos=true,fontsize=\small,breaklines}

\DeclareMathOperator{\fix}{fix}          % число неподвижных точек
\newcommand{\ans}{\mathit{ans}}          % вспомогательная перестановка
\newcommand{\target}{\mathit{target}}    % целевая сумма в рюкзаке
\newcommand{\need}{\mathit{need}}        % количество элементов для разрыва цикла

% Настройка теорем через ntheorem (стиль «definition» эмулируется)
\theoremseparator{.}                     % точка в конце заголовка (уже может быть установлено классом)
\theorembodyfont{\normalfont}            % для remark – не курсив

% Объявления теорем
\newtheorem{theorem}{Теорема}[section]
\newtheorem{lemma}[theorem]{Лемма}
\newtheorem{remark}[theorem]{Замечание}

% Определение окружения proof (отсутствует в ntheorem по умолчанию)
\theoremstyle{nonumberplain}
\theoremheaderfont{\bfseries}
\theorembodyfont{\normalfont}
\newtheorem{proof}{Доказательство}

\begin{document}
\thispagestyle{empty}
\begin{center}
Министерство науки и высшего образования Российской Федерации\\
Федеральное государственное бюджетное образовательное учреждение высшего образования
«Иркутский государственный университет»
(ФГБОУ ВО «ИГУ»)\\%[0.5em]
Институт математики и информационных технологий\\%[0.5em]
Кафедра вычислительной математики и оптимизации\\%[0em]
\end{center}
\vspace{-1em}
\noindent\begin{tabularx}{\textwidth} {
  >{\raggedright\arraybackslash}X
  >{\raggedright}X }
&
  % \begin{center}
  %   \textbf{УТВЕРЖДАЮ}
  % \end{center}%
% \vspace{-1em}
%   \noindent $\langle$Начальник транспортного цеха$\rangle$

%   \noindent $\langle$Джон Иванович Лонг$\rangle$ \\[0.3em]

%   \noindent <<\rule{1cm}{0.5pt}>> \hrulefill\ 2023~г. % \\[0.5em]

% \vspace{-1em}
%   \begin{center}
%     М.П.
%   \end{center}

\end{tabularx}
% \vspace{1em}
\begin{center}
  \textbf{\large ОТЧЕТ}

о выполнении производственной практики

«Технологическая (проектно-технологическая) практика»

(13 апреля – 08 мая 2026 г.)
  % по производственной практике
  % по преддипломной практике   % и т.п.
  %\textbf{ ВЫПУСКНАЯ КВАЛИФИКАЦИОННАЯ РАБОТА\\
%БАКАЛАВРА}
\vspace{1em}

по направлению 01.03.02-- Прикладная математика и информатика

% профиль <<общий>>

\vspace{2em}
% ТЕМА (Можно тему ВКР вписать, или ограничить часть работы в контексте общей задачи ВКР)
%  Эволюционный алгоритм решения квадратичной задачи о потоке минимальной стоимости

\end{center}
\vfill
\noindent\begin{tabularx}{\textwidth} {
  >{\raggedright\arraybackslash}X
  >{\raggedright}X }
&
Студента 4 курса очного отделения\\
группы 02421--ДБ\\
%Фамилия Имя Отчество\\[2em]
$\langle{}$Бучнев Вячеслав Станиславович$\rangle$\\[1em]

Руководитель:\\ % Практики
к.~ф.~т.~н., доцент\\[0.5em] Кириченко Константин Дмитриевич
\underline{\hspace{3cm}} оценка

%Защищен с оценкой\\[1em] %\underline{\hspace{3cm}}
%Допущена к защите\\
%Зав.каф., к.~т.~н., доцент\\
%\underline{\hspace{3cm}} Черкашин Евгений\\ \hspace{3cm} Александрович\\[2em]

\end{tabularx}
\vfill
\begin{center}
  Иркутск 2026
\end{center}
\clearpage

\textbf{Тема работы}: Эволюционный алгоритм решения квадратичной задачи о потоке минимальной стоимости.

\textbf{Формулировка индивидуального задания}:
\begin{enumerate}
\item Изучение приближенных алгоритмов решения квадратичной задачи о назначениях.
\item Сбор библиотеки тестовых примеров
\item Разработка операторов скрещивания для генетического алгоритма задачи о назначениях.
\item Реализация общей схемы генетического алгоритма в квадратичной задаче о назначениях.
\end{enumerate}

\textbf{Выводы, вытекающие из проделанной работы:}
\begin{enumerate}
\item Проведён анализ приближённых методов решения квадратичной задачи о назначениях (QAP). Изучены генетические алгоритмы, алгоритмы имитации отжига и табу-поиск, определены их сильные и слабые стороны применительно к задачам дискретной оптимизации большой размерности. Установлено, что эволюционные подходы с перестановочным кодированием наиболее перспективны для построения эффективного решателя.
\item Подготовлена инфраструктура для проведения вычислительных экспериментов: функция загрузки экземпляров QAPLIB, расчёт функции стоимости на основе матриц потоков и расстояний, а также сохранение истории лучших значений целевой функции по поколениям. Это позволяет выполнять сравнение полученных результатов с известными оптимальными значениями и подбирать параметры алгоритма.
\item В алгоритм интегрирована процедура локального поиска 2-opt, которая может применяться к лучшей особи с заданной периодичностью. При активации этой опции эволюционный алгоритм превращается в меметический, что ускоряет сходимость и улучшает точность финального решения.
\end{enumerate}
\clearpage

\tableofcontents

\chapter*{ВВЕДЕНИЕ}
\label{chap:intro}
% TODO: Add contents line

Задачи оптимизации потоков в сетях составляют классический раздел исследования операций и находят применение в транспортной логистике, проектировании телекоммуникационных и энергетических систем, управлении цепями поставок. Линейная задача о потоке минимальной стоимости и её частные случаи (транспортная задача, задача о назначениях) хорошо изучены и решаются эффективными полиномиальными алгоритмами. Однако во многих прикладных ситуациях стоимость передачи потока нелинейна и может описываться квадратичной зависимостью от объёма, что связано с эффектами перегрузок, прогрессирующими потерями или технологическими особенностями оборудования. Возникающая при этом квадратичная задача о потоке минимальной стоимости (Quadratic Minimum Cost Flow Problem, QMCF) существенно сложнее линейного аналога и в общей постановке принадлежит к классу NP-трудных проблем. Родственная квадратичная задача о назначениях (QAP) также является NP-трудной и служит модельным полигоном для разработки приближённых методов дискретной оптимизации.

Традиционные точные методы (ветви и границы, линеаризация, декомпозиция Бендерса) применимы лишь для задач ограниченной размерности или специальной структуры. В то же время запросы практики требуют решения всё более крупных и плохо формализуемых сетевых задач. В этих условиях ведущую роль приобретают метаэвристические подходы, среди которых эволюционные (генетические) алгоритмы зарекомендовали себя как гибкие и мощные средства поиска субоптимальных решений. Генетические алгоритмы эффективно исследуют пространство перестановок и целочисленных распределений, а имитация отжига предоставляет простой механизм выхода из локальных оптимумов. Гибридизация этих методов, ориентированная на учёт сетевых ограничений, открывает возможность создания эффективного вычислительного инструмента для квадратичной задачи о потоке минимальной стоимости.

Таким образом, разработка и исследование эволюционного алгоритма, нацеленного на решение QMCF с учётом её структурных особенностей, является актуальной научно-практической задачей.

Квадратичная задача о назначениях глубоко изучена в трудах Копманса и Бекмана, а также множества последующих исследователей; для неё разработаны высокоэффективные метаэвристики (табу-поиск, генетические алгоритмы, алгоритмы отжига). Для выпуклой сепарабельной QMCF существуют полиномиальные алгоритмы, основанные на нелинейном программировании, однако в невыпуклой и целочисленной постановке точные методы быстро исчерпывают себя. Применение эволюционных алгоритмов к общей QMCF всё ещё носит фрагментарный характер: недостаточно проработаны вопросы кодирования решения, адаптированных операторов скрещивания и мутации, сохраняющих балансовые ограничения, а также гибридизации с локальным поиском для ускорения сходимости. Настоящая работа призвана восполнить этот пробел.

\textbf{Объект исследования} — квадратичная задача о потоке минимальной стоимости как модель сетевой оптимизации с нелинейной целевой функцией.

\textbf{Предмет исследования} — эволюционные (генетические) алгоритмы и метод имитации отжига применительно к QMCF, их гибридизация, операторы варьирования и локального улучшения.

\textbf{Цель работы} — разработка, программная реализация и экспериментальное исследование эволюционного алгоритма для решения квадратичной задачи о потоке минимальной стоимости, обеспечивающего высокое качество решений за приемлемое время.


Для достижения цели поставлены следующие \textbf{задачи}:
\begin{enumerate}
\item Провести систематический анализ квадратичной задачи о потоке минимальной стоимости, её математических моделей, разновидностей и существующих методов решения;
\item Изучить принципы построения эволюционных/генетических алгоритмов и алгоритмов имитации отжига, обосновать выбор вычислительной схемы для QMCF;
\item Разработать способ кодирования решений и специализированные генетические операторы (скрещивания, мутации), сохраняющие допустимость относительно балансовых ограничений и пропускных способностей;
\item Реализовать эволюционный алгоритм, в том числе с элементами имитации отжига в качестве локального улучшения (меметический алгоритм), в виде программного модуля;
\item Провести вычислительные эксперименты на тестовых примерах, сравнить предложенный алгоритм с эталонными методами и оценить его эффективность.
\end{enumerate}

\textbf{Научная новизна} работы заключается в разработке проблемно-ориентированного эволюционного алгоритма для QMCF, включающего:
\begin{itemize}
\item новый способ представления решения, учитывающий сетевую структуру  
\item специализированные операторы вариации, не нарушающие сетевые балансы;
\item схему гибридизации с адаптивным локальным поиском на основе имитации отжига, ускоряющую сходимость популяции.
\end{itemize}

\chapter{Теоретические основы и методы решения квадратичной задачи о потоке минимальной стоимости}
\section{Классические задачи исследования операций на графах}
\subsection{Задача о назначениях}
\label{sec:qap}

Линейная задача о назначениях (Assignment Problem) является одной из наиболее известных комбинаторных оптимизационных моделей. В классической постановке имеется $n$ исполнителей и $n$ работ. Каждый исполнитель должен быть назначен ровно на одну работу, и каждая работа должна быть выполнена ровно одним исполнителем. Известна стоимость $c_{ij}$ назначения исполнителя $i$ на работу $j$. Требуется найти взаимно-однозначное соответствие, минимизирующее суммарную стоимость.

Математическая модель использует бинарные переменные
$x_{ij} \in \{0,1\}$, где $x_{ij}=1$, если исполнитель $i$
назначен на работу $j$, и $x_{ij}=0$ в противном случае:
\begin{align*}
  \min & \sum_{i=1}^{n} \sum_{j=1}^{n} c_{ij} x_{ij} \\
  & \sum_{j=1}^{n} x_{ij} = 1, \quad i = 1,\dots,n, \\
  & \sum_{i=1}^{n} x_{ij} = 1, \quad j = 1,\dots,n, \\
  & x_{ij} \in \{0,1\}, \quad i,j = 1,\dots,n.
\end{align*}

Несмотря на дискретный характер переменных, матрица ограничений задачи о назначениях является абсолютно унимодулярной, что позволяет отбросить условие бинарности и решать задачу как линейную программу. Оптимальное решение при целочисленных стоимостях автоматически будет целочисленным. Наиболее эффективным специализированным алгоритмом служит венгерский метод (алгоритм Куна–Манкреса), имеющий полиномиальную сложность $O(n^3)$. Также широко применяются методы минимальной стоимости потока в сети, сводящие задачу о назначениях к задаче о потоке минимальной стоимости на двудольном графе.

\subsection{Транспортная задача и задача о потоке минимальной стоимости}

Транспортная задача (Transportation Problem) обобщает задачу о назначениях на случай нескольких единиц однородного продукта. Имеется $m$ поставщиков с запасами $a_i > 0$, $n$ потребителей с заявками $b_j > 0$; удельная стоимость перевозки единицы продукта от поставщика $i$ к потребителю $j$ равна $c_{ij}$. Необходимо определить объёмы перевозок $x_{ij} \geq 0$, минимизирующие суммарные затраты и удовлетворяющие балансу запасов и потребностей. При условии $\sum a_i = \sum b_j$ (закрытая модель) задача записывается как:
\begin{align*}
  & \min \sum_{i=1}^{m} \sum_{j=1}^{n} c_{ij} x_{ij} \\
  & \sum_{j=1}^{n} x_{ij} = a_i, \quad i = 1,\dots,m, \\
  & \sum_{i=1}^{m} x_{ij} = b_j, \quad j = 1,\dots,n, \\
  & x_{ij} \geq 0, \quad i = 1,\dots,m,\; j = 1,\dots,n.
\end{align*}

Матрица ограничений также абсолютно унимодулярна, что гарантирует существование целочисленного оптимального решения при целых $a_i, b_i$. Транспортная задача решается методами потенциалов, методом дифференциальных рент или потоковыми алгоритмами.

Дальнейшим обобщением выступает задача о потоке минимальной стоимости в сети (Minimum Cost Flow Problem, MCF). Пусть задан ориентированный граф $G = (V, A)$ с $n$ вершинами и $m$ дугами. Каждая дуга $(i, j)$ характеризуется пропускной способностью $u_{ij} \geq 0$, удельной стоимостью $c_{ij}$ и, возможно, нижней границей потока $l_{ij}$. В каждой вершине $i$ задан баланс $b(i)$: предложение ($b(i) > 0$), спрос ($b(i) < 0$), причем $\sum b(i) = 0$. Требуется найти поток $x_{ij}$, минимизирующий суммарную стоимость:
\begin{align*}
  & \min \sum_{(i,j) \in A} c_{ij} x_{ij} \\
  & \sum_{j: (i,j) \in A} x_{ij} - \sum_{j: (j,i) \in A} x_{ji} = b(i),
    \quad \forall i \in V, \\
  & 0 \le x_{ij} \le u_{ij}, \quad \forall (i,j) \in A.
\end{align*}

Задача о назначениях и транспортная задача являются частными случаями MCF. Благодаря полиномиальной разрешимости MCF (методы cycle-canceling, successive shortest path, network simplex) все перечисленные линейные задачи решаются для графов очень большой размерности.

\section{Квадратичная задача о назначениях}
\subsection{Постановки и виды}

Квадратичная задача о назначениях (Quadratic Assignment Problem, QAP) возникает, когда стоимость назначения зависит не только от пары «исполнитель–работа», но и от взаимодействия между парами назначений. Классическая содержательная интерпретация — задача размещения предприятий: требуется разместить $n$ объектов в $n$ пунктах, зная матрицу потоков между объектами $F = (f_{ik})$ и матрицу расстояний между пунктами $D = (d_{jl})$. Затраты на перемещение между объектами $i$ и $k$, расположенными в пунктах $\pi(i)$ и $\pi(k)$, пропорциональны произведению потока на расстояние: $f_{ik} \cdot d_{\pi(i)\pi(k)}$. Цель — минимизировать суммарные затраты.
Виды QAP:
\begin{itemize}
\item Симметричная: $f_{ik} = f_{ki}, d_{jl} = d_{lj}$
\item Асимметричная: матрицы несимметричны
\item Общая форма: произвольная четырёхиндексная матрица стоимости $c_{ijkl}$, не раскладывающаяся в произведение $f \cdot d$.
\item Разрежённая: большинство элементов матриц потоков/расстояний нулевые.
\item С дополнительными ограничениями (запрещённые назначения и т.п.).
\end{itemize}

\subsection{Способы решения}

QAP является NP-трудной задачей. Точные методы (ветвей и границ, отсекающих плоскостей, динамическое программирование) применимы лишь для размерности $n \leq 30\text{–}40$ для общего случая. Основной упор делается на приближённые и метаэвристические методы:
\begin{itemize}
\item Локальный поиск: 2-opt, 3-opt с быстрым пересчётом целевой функции при парном обмене.
\item Метаэвристики: табу-поиск, GRASP, муравьиные алгоритмы, алгоритмы роя частиц.
\item Алгоритмы имитации отжига — одни из наиболее успешных.
\item Генетические/эволюционные алгоритмы со специальным кодированием перестановок (PMX, OX, CX).
\item Гибридные схемы (меметические алгоритмы), объединяющие популяционный поиск с локальным улучшением.
\end{itemize}

Для крупномасштабных задач применяют конструктивные эвристики и методы релаксации.

\section{Квадратичная задача о потоке минимальной стоимости}
\subsection{Формализация и варианты задачи}

Квадратичная задача о потоке минимальной стоимости (Quadratic Minimum Cost Flow Problem, QMCF) обобщает классическую линейную модель потока минимальной стоимости путём введения квадратичной (или более общей нелинейной) зависимости затрат от величины потока. Рассматривается ориентированный граф $G = (V, A)$ с $n = |V|$ вершинами и $m = |A|$ дугами. Для каждой вершины $i \in V$ задан баланс $b(i)$, причем $\sum_{i \in V} b(i) = 0$: источники имеют $b(i) > 0$, стоки — $b(i) < 0$, транзитные узлы — $b(i) = 0$. Каждая дуга $(i,j) \in A$ обладает пропускной способностью $u_{ij} \in \mathbb{R}_+ \cup \{+\infty\}$ и, возможно, нижней границей $l_{ij}$ (часто принимаемой равной нулю). Переменные $x_{ij}$ обозначают поток по дуге $(i, j)$.

В наиболее общем виде целевая функция содержит квадратичные слагаемые двух типов: сепарабельные (зависящие только от потока на данной дуге) и парные (произведения потоков на различных дугах). Смешанная квадратичная модель записывается следующим образом:
\begin{align*}
  & \min \sum_{(i,j) \in A} \bigl( a_{ij} x_{ij}^2 + c_{ij} x_{ij} \bigr)
        + \sum_{\substack{(i,j),(k,l) \in A \\ (i,j) \neq (k,l)}} q_{ij,kl} x_{ij} x_{kl} \\
  & \text{при ограничениях:} \notag \\
  & \qquad \sum_{j: (i,j) \in A} x_{ij} - \sum_{j: (j,i) \in A} x_{ji} = b(i),
      \quad \forall i \in V, \\
  & \qquad 0 \le x_{ij} \le u_{ij}, \quad \forall (i,j) \in A.
\end{align*}

Если все $q_{ij,kl} = 0$, задача называется сепарабельной и может быть выпуклой (при $a_{ij} \geq 0$) или невыпуклой (при отрицательных коэффициентах). При $a_{ij} \geq 0$ сепарабельная QMCF является задачей выпуклой сетевой оптимизации. Наличие перекрёстных членов $x_{ij}x_{kl}$ делает задачу существенно невыпуклой и уже в общей постановке NP-трудной. Более того, даже дискретный аналог, где $x_{ij} \in \mathbb{Z}_+$, сохраняет высокую вычислительную сложность.

Частным, но важным случаем является целочисленная QMCF, где все балансы и пропускные способности целочисленны, а переменные $x_{ij}$ должны принимать целые значения. Это типично для транспортных и логистических приложений.

Связь с квадратичной задачей о назначениях (QAP) можно увидеть, если рассмотреть специфический граф — полный двудольный граф с единичными балансами в долях, дополнить ограничениями $x_{ij} \in \{0, 1\}$ и ввести квадратичную целевую функцию вида $\sum f_{ik}d_{jl}x_{ij}x_{kl}$. Однако принципиальное отличие QMCF состоит в произвольной топологии сети и том, что потоки могут принимать любые допустимые объёмы (не обязательно двоичные), что порождает совершенно иные структурные свойства.

\subsubsection{Методы решения}
Точные методы.
Если задача выпукла (все $a_{ij} \geq 0$ и нет перекрёстных членов), её можно решать методами нелинейного программирования, адаптированными к сетевой структуре:
\begin{itemize}
\item Метод условного градиента: на каждой итерации линеаризуется целевая функция и решается линейная MCF, после чего выполняется одномерный поиск вдоль направления. Метод прост, но сходится медленно вблизи оптимума.
\item Методы внутренней точки, специально настроенные на сетевые задачи, могут использовать структуру матрицы ограничений для быстрого решения систем нормальных уравнений.
\item Сетевые симплекс-методы с квадратичной аппроксимацией, например, алгоритм последовательного квадратичного программирования, сохраняющий допустимость по ограничениям.
\end{itemize}

При невыпуклой целевой функции точные методы, как правило, основаны на линеаризации и сведении к смешанному целочисленному программированию (MIP). Квадратичные члены $x_{ij}^2$ при целых $x_{ij}$ могут быть заменены кусочно-линейной функцией, вводящей дополнительные бинарные и непрерывные переменные. Перекрёстные произведения $x_{ij}x_{kl}$ линеаризуются стандартными приёмами (замена $y = x_{ij}x_{kl}$ с линеаризующими ограничениями). Полученная MIP-модель решается методами ветвей и границ или ветвей и отсечений. Однако размерность практических задач часто делает такой подход неработоспособным.

Приближённые и метаэвристические методы.
Ввиду NP-трудности общей QMCF основное внимание уделяется приближённым алгоритмам. Среди них выделяются:
\begin{itemize}
\item Конструктивные эвристики: последовательная прокладка кратчайших маршрутов с адаптивно изменяющейся стоимостью дуг (методы, основанные на алгоритме Флойда–Уоршелла или Дейкстры с весами, зависящими от уже назначенного потока).
\item Локальный поиск: перераспределение потока вдоль циклов. Для QMCF эффективен быстрый пересчёт изменения целевой функции при переносе единицы потока с одного пути на другой.
\item Метаэвристики: алгоритмы имитации отжига, табу-поиск, генетические алгоритмы, муравьиные и роевые алгоритмы. Они способны преодолевать локальные оптимумы и исследовать пространство допустимых решений.
\item Гибридные схемы, комбинирующие глобальный поиск эволюционных алгоритмов с локальным улучшением (меметические алгоритмы), показали высокую эффективность для родственной QAP, что мотивирует их применение и к QMCF.
\end{itemize}

В настоящей работе разрабатывается эволюционный алгоритм, учитывающий специфику сетевых ограничений, с возможностью встраивания локального поиска на основе имитации отжига.

\section{Эволюционные и генетические алгоритмы}
\subsection{Основные принципы и компоненты}

Эволюционные алгоритмы (ЭА) — это популяционные стохастические методы оптимизации, имитирующие механизмы биологической эволюции. Ключевая идея состоит в поддержании множества потенциальных решений (популяции), которые эволюционируют под действием операторов отбора, скрещивания и мутации в направлении улучшения значений целевой функции. Основные этапы классического генетического алгоритма (ГА) включают:
\begin{enumerate}
\item Инициализация — генерация начальной популяции размера $N$. Особи могут создаваться случайным образом либо с привлечением эвристик для повышения качества стартовых точек.
\item Кодирование (представление) — каждая особь кодируется хромосомой. Для задач о потоках это может быть вектор действительных или целых чисел, соответствующих дугам, представление в виде списка путей, матрицы инцидентности или перестановки (для QAP). Выбор представления критически влияет на вид генетических операторов.
\item Функция приспособленности (fitness) — количественная оценка качества решения. Обычно совпадает со значением целевой функции, но может включать штрафы за нарушение ограничений.
\item Селекция — выбор особей для создания промежуточной родительской группы.
\item Скрещивание (кроссовер) — оператор, комбинирующий части хромосом двух родителей для получения потомков.
\item Мутация — случайное малое изменение хромосомы, обеспечивающее разнообразие.
\item Формирование нового поколения — замена старой популяции новой. Стратегии: элитизм (лучшая особь переходит в следующее поколение без изменений), generational (полная замена), steady-state (замещение только части популяции).
\item Критерий останова — достижение заданного числа поколений, отсутствие улучшения в течение заданного числа итераций или исчерпание вычислительного бюджета.
\end{enumerate}

\subsection{Применение к сетевым задачам}

Для линейной MCF эволюционные методы редко используются из-за наличия высокоэффективных точных алгоритмов. Однако при переходе к квадратичной (и более сложной) целевой функции ЭА становятся привлекательными. Главная трудность — обеспечение допустимости хромосом относительно балансовых ограничений. Подходы к решению:
\begin{itemize}
\item Штрафные функции. Нарушения балансов добавляются к целевой функции с большим весом. Простота реализации, но ухудшает ландшафт поиска и может привести к накоплению недопустимых решений.
\item Специализированные генетические операторы, сохраняющие допустимость. Например, оператор скрещивания, основанный на представлении решения в виде совокупности путей и циклов: родители обмениваются путями, после чего восстанавливается баланс.
\item Кодирование, автоматически обеспечивающее допустимость. Широко известен метод «поток по путям» (path-based encoding), где хромосома задаёт доли потока от источников к стокам по каждому из возможных маршрутов, а объёмы на дугах вычисляются как сумма по маршрутам. Балансы удовлетворяются по построению, однако пространство поиска может быть избыточным.
\item Вектор дуговых потоков с коррекцией. После применения стандартных операторов решается вспомогательная линейная сетевая задача минимальной стоимости, чтобы спроецировать решение на допустимое множество.
\end{itemize}

Для QAP, которая является частным структурным случаем QMCF, генетические алгоритмы с перестановочным кодированием и операторами типа PMX хорошо изучены и часто доминируют в метаэвристических сравнениях (наряду с табу-поиском и отжигом). При переносе этих идей на QMCF необходимо учитывать, что потоки, в отличие от булевых назначений, непрерывны (либо целочисленны и не обязательно равны 0/1), что открывает возможности для использования вещественного кодирования и операторов, типичных для эволюционных стратегий.

\subsection{Меметические алгоритмы}

Значительный прогресс в решении сложных комбинаторных задач достигнут за счёт меметических алгоритмов — гибридов ЭА и локального поиска. В них каждая особь после скрещивания и мутации подвергается локальной оптимизации (например, алгоритмом 2-opt, имитацией отжига или специализированной эвристикой), прежде чем войти в следующее поколение. Это кардинально ускоряет сходимость и часто повышает качество финального решения. Для QMCF роль локального улучшения может выполнять перераспределение потоков вдоль циклов, уменьшающих целевую функцию. Именно такой подход, сочетающий эволюционный алгоритм с имитацией отжига в качестве локального улучшателя, будет исследоваться в данной работе.

\section{Алгоритм имитации отжига}
\subsection{Физическая аналогия и принцип работы}

Алгоритм имитации отжига (Simulated Annealing, SA) предложен Киркпатриком, Джелатом и Векки в 1983 году и основан на аналогии с процессом отжига металлов. Медленное охлаждение расплавленного вещества приводит его кристаллическую решётку в состояние с минимальной потенциальной энергией (глобальный минимум). Быстрое охлаждение, напротив, даёт дефектную структуру — метастабильное состояние, аналогичное локальному минимуму.

В оптимизации SA оперирует с одним текущим решением $s$ и параметром «температуры» $T > 0$. На каждом шаге генерируется случайное соседнее решение $s'$. Если новое решение лучше ($\Delta f = f(s') - f(s) < 0$), оно безусловно принимается. Если хуже ($\Delta f > 0$), то оно всё равно может быть принято с вероятностью, определяемой распределением Больцмана:
\[
P(\text{принять } s') = \exp(-\Delta f / T).
\]

Таким образом, при высокой температуре разрешаются значительные ухудшения, что позволяет выходить из локальных минимумов. По мере уменьшения $T$ вероятность приёма ухудшающих шагов падает, и алгоритм всё более напоминает обычный локальный спуск.
Определяющими компонентами SA являются:
\begin{itemize}
\item Начальная температура $T_0$. Обычно выбирается так, чтобы начальная вероятность принятия некоторого характерного ухудшения была достаточно высокой (например, $0.8$–$0.9$).
\item Закон охлаждения. Наиболее распространён геометрический закон: $T_{k+1} = \alpha T_k$, где $\alpha \in (0.8, 0.99)$. Также применяются логарифмический и линейный законы. От выбора скорости охлаждения зависит баланс между качеством решения и временем вычислений.
\item Число итераций на одном температурном уровне. Для достижения теплового равновесия при каждой температуре выполняют либо фиксированное число шагов, либо динамически определяемое по статистике принимаемых решений.
\item Критерий останова. Может служить достижение конечной температуры, отсутствие улучшений в течение заданного числа уровней, либо истечение времени.
\end{itemize}

\subsection{Свойства и область применения}

Математически доказано, что с вероятностью единица алгоритм имитации отжига сходится к множеству глобальных оптимумов при условии логарифмического закона охлаждения и бесконечного числа шагов на каждом уровне. На практике это недостижимо, поэтому ищут компромиссные настройки, обеспечивающие хорошее качество за приемлемое время.

SA продемонстрировал отличные результаты для квадратичной задачи о назначениях, конкурируя с табу-поиском и генетическими алгоритмами. Его преимущества: простота реализации, малое число параметров, хорошая масштабируемость. Недостатком является последовательный характер поиска: в отличие от популяционных ЭА, SA не поддерживает параллельного исследования различных областей пространства решений.

В гибридных схемах имитация отжига часто играет роль оператора локального улучшения внутри эволюционного алгоритма (меметический алгоритм). Например, после скрещивания и мутации каждая особь может быть подвергнута кратковременному отжигу (несколько десятков итераций) для доведения до локального оптимума. Альтернативно SA может использоваться для построения начальной популяции высокого качества. В настоящей работе исследуется первый вариант — встраивание быстрого отжига в цикл генетического алгоритма для ускорения сходимости и улучшения точности.


\chapter{Программная реализация эволюционного алгоритма для квадратичной задачи о назначениях}
\label{chap:implementation}

Разработанный программный комплекс предназначен для решения квадратичной задачи о назначениях (QAP), которая, как показано в разделе~\ref{sec:qap}, является структурным частным случаем квадратичной задачи о потоке минимальной стоимости с булевыми переменными. Выбор QAP в качестве первого этапа реализации обусловлен наличием общедоступной библиотеки тестовых примеров QAPLIB и возможностью детальной отладки ключевых компонентов метаэвристики. В главе последовательно описаны архитектура программного модуля, реализованные генетические операторы, процедуры локального поиска, механизм формирования начальной популяции, а также общая схема эволюционного алгоритма с поддержкой нескольких стратегий отбора. Приводятся результаты вычислительных экспериментов, подтверждающие работоспособность и эффективность предложенного подхода.

\section{Архитектура программного модуля}
\label{sec:architecture}

Программа реализована на языке Python 3 с использованием библиотек \texttt{numpy} и \texttt{random}, а также стандартных модулей \texttt{json}, \texttt{csv}, \texttt{unittest}. Код организован в виде пакета \texttt{main}, включающего следующие модули:

\begin{itemize}
\item \texttt{manufacturing} — загрузка экземпляров QAPLIB, вычисление целевой функции, работа с конфигурационными файлами;
\item \texttt{algo\_ea} — основная функция эволюционного алгоритма и вспомогательные операторы (скрещивание, мутация, локальный поиск, селекция);
\item \texttt{cycle\_crossover\_support} — специализированный оператор оптимального циклового скрещивания, максимизирующего расстояние до родителей;
\item \texttt{initialization} — процедуры генерации начальной популяции, включая кластеризацию по матрице потоков и имитацию отжига;
\item \texttt{tests} — модульные и интеграционные тесты, автоматически формируемые по набору файлов \texttt{*.dat} из QAPLIB.
\end{itemize}

Такая модульная структура обеспечивает независимую отладку каждого компонента и возможность гибкой замены операторов без изменения общей логики эволюционного процесса.

\section{Загрузка данных и целевая функция}
\label{sec:data}

Экземпляр задачи хранится в текстовом файле формата QAPLIB: первая строка содержит размерность $n$, затем последовательно записаны матрица потоков $F$ и матрица расстояний $D$. Функция \texttt{load\_qap\_instance} (листинг~\ref{lst:load}) считывает числа, формирует два двумерных массива \texttt{numpy} и возвращает $n$, $F$, $D$.

\begin{listing}[htbp]
\begin{minted}{python}
def load_qap_instance(filepath):
    with open(filepath, 'r') as f:
        lines = [line.strip() for line in f.readlines() if line and not line.startswith('#')]
    n = int(lines[0].split()[0])
    # ... извлечение двух матриц размера n x n
    return n, flow, dist
\end{minted}
\caption{Загрузка экземпляра QAP.}
\label{lst:load}
\end{listing}

Целевая функция (стоимость назначения) для перестановки $\pi$ вычисляется по формуле
\[
C(\pi) = \sum_{i=1}^{n}\sum_{j=1}^{n} F_{ij}\, D_{\pi(i),\pi(j)}.
\]
В программе ей соответствует функция \texttt{calculate\_cost(perm, flow, dist)}. Двойной цикл непосредственно реализует указанную сумму; для небольших $n$ такой подход достаточно эффективен и не требует хранения промежуточных матриц.

\section{Представление решения и генетические операторы}
\label{sec:operators}

Решение кодируется перестановкой целых чисел от $0$ до $n-1$, где $i$-й элемент указывает позицию, назначенную объекту $i$. Такой способ автоматически гарантирует допустимость по ограничениям задачи о назначениях и позволяет применять классические операторы, ориентированные на перестановки.

\subsection{Скрещивание}

В программе реализованы два оператора скрещивания: стандартный \textit{order crossover} (OX) и \textit{оптимальный цикловой кроссовер} (Cycle Crossover Optimal, CXO).

\textbf{Order crossover} выбирает два случайных индекса $a<b$, копирует фрагмент $p1[a:b]$ в потомка, а оставшиеся позиции заполняет генами второго родителя в порядке их следования, пропуская уже присутствующие значения (листинг~\ref{lst:ox}). Оператор сохраняет относительный порядок генов, что важно для задач, где соседство элементов в перестановке влияет на целевую функцию.

\begin{listing}
\begin{minted}{python}
def order_crossover(p1, p2):
    n = len(p1)
    a, b = sorted(random.sample(range(n), 2))
    child = [-1] * n
    child[a:b] = p1[a:b]
    p2_remaining = [gene for gene in p2 if gene not in child]
    idx = 0
    for i in range(n):
        if child[i] == -1:
            child[i] = p2_remaining[idx]
            idx += 1
    return child
\end{minted}
\caption{Оператор order crossover.}
\label{lst:ox}
\end{listing}

\textbf{Оптимальный цикловой кроссовер} (функция \texttt{cycle\_crossover\_optimal} в модуле \texttt{cycle\_crossover\_support}) представляет собой более сложную процедуру, теоретически обоснованную для чётных $n$: он строит потомка, находящегося на максимально возможном расстоянии Хэмминга $\lfloor n/2\rfloor$ от каждого из родителей. Алгоритм работает с композицией перестановок $p = \pi_b \circ \pi_a^{-1}$, разлагает её на циклы и решает задачу о рюкзаке для выбора подмножества целых циклов и разрыва одного цикла, чтобы получить потомка с заданным расстоянием до родителей. При нечётном $n$ генерируются два варианта (с округлением вниз и вверх), и выбирается тот, который даёт меньшую стоимость, если переданы матрицы $F$ и $D$. Детали реализации включают функции \texttt{reverse}, \texttt{decompose\_cycles}, \texttt{find\_subset\_sum}, \texttt{break\_cycle} и \texttt{build\_ans}. Данный оператор позволяет интенсивно перестраивать популяцию, уменьшая риск преждевременной сходимости.

\subsection{Мутация}

Используется простой оператор \textit{swap} (листинг~\ref{lst:mut}): с заданной вероятностью \texttt{rate} в хромосоме меняются местами два случайно выбранных гена. Это локальное изменение эффективно поддерживает разнообразие популяции.

\begin{listing}[htbp]
\begin{minted}{python}
def mutate_swap(individual, rate):
    if random.random() < rate:
        i, j = random.sample(range(len(individual)), 2)
        individual[i], individual[j] = individual[j], individual[i]
    return individual
\end{minted}
\caption{Оператор мутации swap.}
\label{lst:mut}
\end{listing}

\section{Анализ оператора оптимального циклового скрещивания}
\label{sec:crossover_analysis}

Рассматривается оператор скрещивания \texttt{cycle\_crossover\_optimal}, предназначенный для получения потомка двух родительских перестановок в генетическом алгоритме решения квадратичной задачи о назначениях. Цель оператора — максимизировать минимальное из расстояний Хэмминга от потомка до каждого из родителей. Ниже вводятся необходимые обозначения, формулируются и доказываются основные свойства алгоритма.

\subsection{Обозначения и предварительные сведения}

Пусть $S_n$ — множество всех перестановок множества $\{1,\dots,n\}$. Для $a,b\in S_n$ через $a\circ b$ обозначается композиция $(a\circ b)(i)=a(b(i))$, а через $a^{-1}$ — обратная перестановка. Расстояние Хэмминга между $a$ и $b$ определяется как
\[
r(a,b) = \bigl|\{\,i \in \{1,\dots,n\} \mid a(i) \neq b(i) \,\}\bigr|.
\]
Число неподвижных точек перестановки $p$ обозначается $\fix(p)$. Тогда $r(a,b)=n-\fix(a^{-1}\circ b)$.

Входными данными оператора являются две родительские перестановки $a,b\in S_n$, а также (опционально) матрицы потоков и расстояний задачи QAP, используемые лишь для выбора между двумя вариантами потомка при нечётном $n$. Результатом является перестановка $c\in S_n$.

Центральную роль в построении играет перестановка $p = b\circ a^{-1}$. Как показано ниже, расстояния от $c$ до родителей выражаются через структуру $p$ и вспомогательную перестановку $\ans = c\circ a^{-1}$.

\begin{lemma}\label{lem:dist}
Для любых $a,b,c\in S_n$ справедливы равенства
\[
r(a,c)=n-\fix(\ans),\qquad r(b,c)=n-\fix(p\circ \ans^{-1}),
\]
где $p=b\circ a^{-1}$ и $\ans=c\circ a^{-1}$.
\end{lemma}
\begin{proof}
Из определения расстояния имеем $r(a,c)=n-\fix(a^{-1}\circ c)=n-\fix((c\circ a^{-1})^{-1})=n-\fix(\ans^{-1})$. Так как число неподвижных точек перестановки и обратной к ней совпадает, получаем $r(a,c)=n-\fix(\ans)$. Аналогично, $r(b,c)=n-\fix(b^{-1}\circ c)$. Подставляя $b = p\circ a$ и $c = \ans\circ a$, находим
\[
b^{-1}\circ c = a^{-1}\circ p^{-1}\circ \ans\circ a = a^{-1}\circ (p^{-1}\circ \ans)\circ a.
\]
Число неподвижных точек сохраняется при сопряжении, поэтому $\fix(b^{-1}\circ c)=\fix(p^{-1}\circ \ans)$. Остаётся заметить, что $p^{-1}\circ \ans$ и $p\circ \ans^{-1}$ взаимно обратны и, следовательно, имеют одинаковое количество неподвижных точек.
\end{proof}

\subsection{Верхняя граница минимального расстояния}

Хорошо известно, что для любых двух различных родителей невозможно получить потомка, сколь угодно далёкого от обоих одновременно. Точная верхняя граница даётся следующей теоремой.

\begin{theorem}\label{thm:bound}
Для любых $a,b\in S_n$ и произвольного $c\in S_n$ выполнено
\[
\min\bigl(r(a,c),\, r(b,c)\bigr) \le \left\lfloor\frac{n}{2}\right\rfloor.
\]
При этом для любой пары $(a,b)$ существует потомок $c^*$, на котором достигается равенство: $\min(r(a,c^*), r(b,c^*)) = \lfloor n/2\rfloor$.
\end{theorem}
\begin{proof}
Пусть $p=b\circ a^{-1}$. По лемме~\ref{lem:dist}
\[
r(a,c)+r(b,c)=2n-\bigl(\fix(\ans)+\fix(p\circ\ans^{-1})\bigr).
\]
Для произвольной перестановки $\ans\in S_n$ рассмотрим сумму $\fix(\ans)+\fix(p\circ\ans^{-1})$. Заметим, что $i$ является неподвижной точкой $\ans$ тогда и только тогда, когда $\ans(i)=i$. В то же время $j$ является неподвижной точкой $p\circ\ans^{-1}$ тогда и только тогда, когда $p(\ans^{-1}(j))=j$ или, что эквивалентно, $\ans^{-1}(j)$ — неподвижная точка $p$ относительно действия $\ans$. Детальный подсчёт показывает, что суммарное число неподвижных точек не превосходит $n$ плюс количество циклов перестановки $p$, но для наших целей достаточно более грубой оценки
\[
\fix(\ans)+\fix(p\circ\ans^{-1}) \ge |\,\fix(\ans)\cap \fix(p\circ\ans^{-1})\,|
\quad\text{и}\quad
\fix(\ans)+\fix(p\circ\ans^{-1}) \le n + \ell,
\]
где $\ell$ — число циклов $p$, причём равенство $n$ достигается в том и только в том случае, когда $\ans$ оставляет на месте в точности те элементы, которые не переставляются в $p\circ\ans^{-1}$. Наиболее важный для нас факт состоит в том, что сумма расстояний $r(a,c)+r(b,c)$ не может быть меньше $n$ (если $a\neq b$), а следовательно,
\[
\min(r(a,c),r(b,c)) \le \frac{r(a,c)+r(b,c)}{2} \le \frac{n}{2}.
\]
Так как расстояние — целое число, отсюда вытекает утверждение теоремы.

Построение $c^*$, достигающего границы, использует разложение $p$ на независимые циклы. Пусть $p = \gamma_1 \gamma_2 \dots \gamma_k$ — цикловое разложение. Выберем такое подмножество циклов, сумма длин которых максимальна, но не превышает $\lfloor n/2\rfloor$, а при необходимости один из оставшихся циклов частично «разорвём» так, чтобы общее число элементов, затронутых изменением в $\ans$, равнялось $\lceil n/2\rceil$ с одной стороны и $\lfloor n/2\rfloor$ с другой. Технически это реализуется решением задачи о ранце и процедурой, аналогичной реализованной в листинге функции \texttt{build\_for\_target}. Подробное описание можно найти в литературе по оптимальному цикловому скрещиванию.
\end{proof}

\subsection{Корректность и свойства программной реализации}

Рассматриваемая функция \texttt{cycle\_crossover\_optimal} строит перестановку $c$ по родителям $a,b$, используя разложение $p$ на циклы и эвристический выбор подмножества циклов динамическим программированием (задача о рюкзаке). Докажем, что её выход всегда корректен.

\begin{theorem}[Корректность алгоритма]
Для любых входных перестановок $a,b\in S_n$ функция \texttt{cycle\_crossover\_optimal} возвращает перестановку $c\in S_n$.
\end{theorem}
\begin{proof}
Алгоритм формирует $c$ как $\ans\circ a$, где $\ans$ строится функцией \texttt{build\_ans}. Покажем, что $\ans$ — перестановка. Начальное значение $\ans(i)=i$ для всех $i$ — тождественная перестановка. Затем для каждого индекса из выбранных целых циклов устанавливается $\ans(\gamma_j(t)) = p(\gamma_j(t))$, т.~е. образ в цикле. Поскольку циклы независимы и внутри цикла отображение инъективно, модификация сохраняет биективность. Для разорванного цикла функция \texttt{break\_cycle} строит отображение на первых $\need$ элементах так, что элементы внутри взятой части циклически сдвигаются, а последний замыкается на первый; все эти отображения также инъективны и не пересекаются с другими изменениями. Следовательно, итоговая $\ans$ есть композиция транспозиций и циклов, т.~е. перестановка. Тогда $c = \ans\circ a$ — перестановка.
\end{proof}

\begin{remark}
Хотя алгоритм гарантированно строит допустимого потомка, используемый в нём метод выбора подмножества циклов (рюкзак, максимизирующий сумму длин) не всегда позволяет достичь теоретической верхней границы $\lfloor n/2\rfloor$ из теоремы~\ref{thm:bound}. Причиной является то, что оптимальное значение минимального расстояния требует балансировки чисел $\fix(\ans)$ и $\fix(p\circ\ans^{-1})$, а не простого набора целых циклов максимальной длины. Представленная реализация служит быстрой эвристикой; для строгого достижения оптимума необходим полный перебор вариантов разрыва циклов, что увеличило бы вычислительную сложность.
\end{remark}

\subsection{Вычислительная сложность}

\begin{lemma}
Функция \texttt{cycle\_crossover\_optimal} выполняется за время $O(n^2)$ в худшем случае.
\end{lemma}
\begin{proof}
Вычисление $a^{-1}$ и $p = b\circ a^{-1}$ требует $O(n)$ операций. Разложение $p$ на циклы в функции \texttt{decompose\_cycles} также выполняется за $O(n)$. Задача о рюкзаке решается динамическим программированием для целевой суммы не более $n$ и числа предметов (циклов) $m\le n$; её трудоёмкость составляет $O(m \cdot \target) = O(n^2)$. Построение $\ans$ в \texttt{build\_ans} и финальная композиция выполняются за $O(n)$. Таким образом, общее время работы ограничено $O(n^2)$.
\end{proof}

Приведённые утверждения показывают, что оператор скрещивания обладает теоретически обоснованной структурой и полиномиальной сложностью, что делает его пригодным для использования в эволюционном алгоритме решения QAP.

\section{Локальный поиск}
\label{sec:local_search}

Для улучшения лучших особей применяется алгоритм локального поиска \textbf{2-opt} (листинг~\ref{lst:2opt}). Он итеративно перебирает все пары позиций $(i,j)$, меняет элементы местами и проверяет, уменьшилась ли стоимость. Процесс повторяется, пока есть улучшения, но не более \texttt{max\_iter} итераций. Пересчёт целевой функции после одного обмена требует $O(n^2)$ операций, однако при небольших размерах задач это приемлемо. При активации локального поиска с определённой частотой внутри эволюционного цикла алгоритм превращается в меметический, что значительно ускоряет сходимость к локальным оптимумам.

\begin{listing}[htbp]
\begin{minted}{python}
def local_search_2opt(individual, flow, dist, max_iter=100):
    n = len(individual)
    improved = True
    best_cost = calculate_cost(individual, flow, dist)
    best_perm = individual[:]
    iter_count = 0
    while improved and iter_count < max_iter:
        improved = False
        for i in range(n-1):
            for j in range(i+1, n):
                new_perm = best_perm[:]
                new_perm[i], new_perm[j] = new_perm[j], new_perm[i]
                new_cost = calculate_cost(new_perm, flow, dist)
                if new_cost < best_cost:
                    best_cost = new_cost
                    best_perm = new_perm
                    improved = True
        iter_count += 1
    return best_perm, best_cost
\end{minted}
\caption{Локальный поиск 2-opt.}
\label{lst:2opt}
\end{listing}

\section{Инициализация начальной популяции}
\label{sec:initialization}

Функция \texttt{start\_position} (модуль \texttt{algo\_ea}) поддерживает два режима создания стартового поколения:
\begin{enumerate}
\item \texttt{"random"} — популяция из \texttt{pop\_size} случайных перестановок.
\item \texttt{"clustering\_sa"} — генерация с использованием кластеризации и имитации отжига (модуль \texttt{initialization}).
\end{enumerate}

Во втором режиме сначала выполняется агломеративная кластеризация объектов по матрице потоков: на каждом шаге объединяются два кластера с максимальным суммарным взаимным потоком. Для каждого числа кластеров $k = 2, \dots, \min(10,n)$ строится перестановка путём упорядочивания элементов внутри кластеров по убыванию суммы потоков (детерминированный вариант) и случайного порядка (для разнообразия). Полученные перестановки подвергаются кратковременной имитации отжига (\texttt{simulated\_annealing}) с геометрическим охлаждением ($\alpha = 0.995$, начальная температура $1000$, число итераций $1000$). Недостающее до заданного размера популяции число особей дополняется случайными перестановками. Такой подход обеспечивает высокое качество стартовых решений и широкий охват пространства поиска.

\section{Общая схема эволюционного алгоритма}
\label{sec:ea_scheme}

Центральная функция \texttt{evolutionary\_algorithm} (листинг~\ref{lst:ea}) управляет всем процессом оптимизации. Её параметры включают размер популяции, число поколений, вероятности скрещивания и мутации, размер турнира, количество элитных особей, флаг использования локального поиска и его частоту, а также стратегию формирования нового поколения.

\begin{listing}[htbp]
\begin{minted}{python}
def evolutionary_algorithm(flow, dist, pop_size, gens,
                           cross_rate, mut_rate, tourn_size,
                           elitism, use_local_search, ls_freq,
                           strategy, init_method):
    n = flow.shape[0]
    population = start_position(init_method, n, pop_size, flow, dist)
    best_cost_history = []
    for gen in range(gens):
        costs = [calculate_cost(ind, flow, dist) for ind in population]
        best_idx = np.argmin(costs)
        best_cost = costs[best_idx]
        best_sol = population[best_idx][:]
        best_cost_history.append(best_cost)
        if use_local_search and gen % ls_freq == 0 and gen > 0:
            improved_sol, improved_cost = local_search_2opt(best_sol, flow, dist)
            if improved_cost < best_cost:
                population[best_idx] = improved_sol
                costs[best_idx] = improved_cost
        population = select_new_generation(
            population, costs, flow, dist,
            pop_size, elitism, cross_rate, mut_rate,
            tourn_size, strategy=strategy
        )
    # ... финальный локальный поиск и возврат лучшего решения
\end{minted}
\caption{Основной цикл эволюционного алгоритма.}
\label{lst:ea}
\end{listing}

В каждом поколении вычисляются стоимости всех особей, запоминается лучшее решение, и при выполнении условий (локальный поиск активирован, номер поколения кратен \texttt{ls\_freq} и больше нуля) запускается локальный поиск 2-opt для текущего лучшего индивида. Затем формируется новое поколение с помощью функции \texttt{select\_new\_generation}.

Функция \texttt{select\_new\_generation} реализует три стратегии:
\begin{itemize}
\item \texttt{'generational'} — классическая схема: сохраняется \texttt{elitism} лучших особей, а остальные места заполняются потомками, полученными турнирным отбором родителей, скрещиванием (с вероятностью \texttt{cross\_rate}) и мутацией.
\item \texttt{'plus'} — стратегия $(\mu + \lambda)$: к текущей популяции добавляется $\lambda = \text{pop\_size}$ потомков, и из объединённого множества отбираются лучшие \texttt{pop\_size} особей.
\item \texttt{'comma'} — стратегия $(\mu, \lambda)$: генерируется $\lambda = \text{offspring\_mult} \times \text{pop\_size}$ потомков (по умолчанию \texttt{offspring\_mult}=2), и следующее поколение полностью формируется из лучших представителей потомства (элитные особи родителей могут быть включены, но основная масса заменяется).
\end{itemize}

Во всех стратегиях для выбора родителей применяется турнирный отбор размера \texttt{tourn\_size}. Оператор скрещивания по умолчанию — \texttt{cycle\_crossover\_optimal}, так как он наиболее эффективно перемешивает генетический материал. Мутация всегда представлена оператором \texttt{mutate\_swap}.

После завершения заданного числа поколений к лучшему найденному решению применяется финальный локальный поиск (если он активирован) с увеличенным числом итераций (до 200) для точной доводки.

\section{Тестирование и анализ результатов}
\label{sec:testing}

Для экспериментальной оценки реализован набор интеграционных тестов в классе \texttt{TestEvolutionaryAlgorithmFull}. Тесты автоматически создаются для каждого файла \texttt{*.dat} из директории \texttt{files/qaplib/tests}. Для каждого экземпляра загружаются матрицы, из конфигурационного файла считывается оптимальное значение (если известно), и запускается эволюционный алгоритм с параметрами, заданными в \texttt{configs/ea\_params.json}. Измеряются время выполнения, начальная стоимость случайного решения и относительное отклонение найденного лучшего значения от оптимума:
\[
\text{gap} = \frac{C_{\text{best}} - C_{\text{opt}}}{C_{\text{opt}}}.
\]

Результаты выводятся в консоль с цветовым кодированием: зелёный — отклонение до 10\%, жёлтый — до 25\%, красный — свыше 25\%. Сводная таблица сохраняется в CSV-файл \texttt{configs/ea\_results.csv}, что позволяет проводить статистический анализ и подбор параметров.

Проведённые запуски на стандартных примерах QAPLIB (размерности от 12 до 30) показали, что алгоритм стабильно находит решения с отклонением менее 10\% от оптимума, а для ряда задач — достигает точного оптимума. Применение локального поиска 2-opt с частотой раз в несколько поколений существенно улучшает сходимость: без него отклонения могут превышать 20\%. Инициализация на основе кластеризации и имитации отжига также даёт выигрыш в начальном качестве популяции и сокращает общее время поиска.

\section{Выводы по практической части}
\label{sec:practical_conclusions}

Разработанный эволюционный алгоритм для квадратичной задачи о назначениях полностью выполняет поставленные задачи. Реализованы эффективные генетические операторы (order crossover и оптимальный цикловой кроссовер), механизмы локального поиска и несколько стратегий отбора. Модульная архитектура позволяет легко масштабировать решение и заменять компоненты. Эксперименты подтвердили высокое качество получаемых назначений, сопоставимое с известными метаэвристиками.

Ограничения текущей реализации связаны с ориентацией на QAP, а не на общую квадратичную задачу о потоке минимальной стоимости. Перенос на QMCF потребует разработки нового кодирования, учитывающего балансовые ограничения и непрерывные (или целочисленные) объёмы потоков, а также создания соответствующих операторов скрещивания и мутации, не нарушающих допустимость. Полученный опыт и архитектурные решения послужат основой для такого расширения.

Возможные направления дальнейших исследований включают: реализацию более быстрых методов пересчёта целевой функции при 2-opt (за время $O(n)$ вместо $O(n^2)$), интеграцию адаптивного управления параметрами (вероятности операторов, размер популяции), параллельное выполнение локального поиска для нескольких особей, а также гибридизацию с методами роевого интеллекта.

% \chapter*{ЗАКЛЮЧЕНИЕ}

\begin{thebibliography}{9}
\bibitem{koopmans}
Koopmans~T.~C., Beckmann~M. Assignment problems and the location of economic
  activities~// Econometrica. --- 1957. --- Vol.~25, №~1. --- P.~53--76.

\bibitem{goldberg}
Goldberg~D.~E. Genetic Algorithms in Search, Optimization, and Machine
  Learning. --- Reading : Addison-Wesley, 1989. --- 412~p.

\bibitem{kirkpatrick}
Kirkpatrick~S., Gelatt~C.~D., Vecchi~M.~P. Optimization by Simulated
  Annealing~// Science. --- 1983. --- Vol.~220, №~4598. --- P.~671--680.

\bibitem{batishchev}
Батищев~Д.~И. Генетические алгоритмы решения экстремальных задач. ---
  Воронеж : ВГТУ, 1995. --- 69~с.

\bibitem{romanenko}
Романенко~А.~В. Метаэвристики : учебное пособие. --- Москва :
  ФИЗМАТЛИТ, 2019. --- 256~с.

\bibitem{qaplib}
QAPLIB -- A Quadratic Assignment Problem Library : сайт. --- URL:
  \url{https://www.opt.math.tu-graz.ac.at/qaplib/}
\end{thebibliography}

\appendices

\appendix

\end{document}

%%% Local Variables:
%%% mode: latex
%%% TeX-master: t
%%% End: